\chapter{Peluang dan Peluang Bersyarat}

Sebuah klinik memasang iklan uji saring untuk penyakit yang jarang: hanya satu dari
seribu orang di daerah itu mengidapnya. Ujinya memang bagus. Dari setiap seratus orang
sakit, sembilan puluh sembilan memberi hasil positif. Dari setiap seratus orang sehat,
hanya satu yang salah dituduh positif. Iklan menuliskan kedua angka itu sebagai
``ketepatan 99 persen''.

Sekarang hasil positif itu ada di tangan kita. Berapa peluang kita benar-benar sakit?
Pada populasi $100{,}000$ orang, sekitar $100$ mengidap penyakit itu dan $99$ di
antaranya positif. Sisanya, $99{,}900$ orang sehat, menyumbang $999$ positif palsu.
Jadi ada $1{,}098$ hasil positif dan hanya $99$ yang benar: $99/1{,}098=11/122$,
yaitu $0.0902$. Sekitar satu dari sebelas, bukan sembilan puluh sembilan dari seratus.

Kedua angka itu benar dan keduanya menjawab pertanyaan yang berbeda. Iklan memberi
peluang hasil positif bila orangnya sakit. Pasien ingin tahu peluang orangnya sakit
bila hasilnya positif. Menukar dua arah itu bukan kesalahan aritmetika; itu kesalahan
menamai besaran, dan akibatnya sering sepuluh kali lipat.

Bab 2 sudah memberi bahasa untuk kejadian: gabungan, irisan, komplemen. Bab 5 memberi
cara menjumlahkan tak hingga banyak suku. Keduanya belum cukup, karena himpunan tidak
membawa bobot: $A\cup B$ tetap himpunan yang sama, entah $A$ hampir pasti terjadi atau
hampir pasti tidak. Yang dibutuhkan tiga hal: fungsi bobot dengan aturan yang tidak boleh
dilanggar, cara mencacah bobot, dan aturan memperbarui bobot ketika informasi baru masuk.
Bab ini menyusun ketiganya, lalu memeriksanya dengan simulasi berbenih tetap.

\begin{untukanak}
Bayangkan penjaga gerbang sekolah memeriksa tas dengan alat yang berbunyi. Alat itu
bagus: dari seratus tas yang benar-benar membawa barang terlarang, sembilan puluh
sembilan membuatnya berbunyi. Tetapi alat itu juga berbunyi untuk satu dari seratus
tas yang isinya biasa saja. Sekarang misalkan di sekolah itu hanya satu dari seribu
tas yang benar-benar membawa barang terlarang.

Pertanyaannya sederhana dan mengejutkan. Kalau alat berbunyi, apakah tas itu pasti
berbahaya? Ternyata tidak. Tas berbahaya memang jarang, sedangkan tas biasa sangat
banyak, sehingga bunyi palsu dari tas biasa jauh lebih banyak daripada bunyi benar.
Penjaga yang hanya melihat lampu merah akan menahan banyak anak yang tidak bersalah.

Pertanyaan yang membuat penasaran: berapa bagian dari semua bunyi itu yang benar?
Jawabannya dapat dihitung tanpa menebak, asalkan kita tahu berapa banyak tas dari
masing-masing jenis. Di bab ini kita akan tahu mengapa bunyi alat yang ``sembilan puluh
sembilan persen akurat'' bisa berubah menjadi tuduhan yang salah dalam sembilan dari
sepuluh kasus, dan mengapa arah pertanyaannya menentukan jawabannya.
\end{untukanak}

\section{Aksioma peluang dan akibatnya}
\tingkat{dasar}

Sebelum satu angka pun dihitung, kita harus menyepakati benda apa yang diberi bobot.
Kalimat sehari-harinya: peluang adalah bobot yang kita taruh pada setiap kemungkinan,
dengan aturan bahwa total bobot seluruh kemungkinan adalah satu, dan bobot tidak pernah
negatif. Aturan itu tidak menjelaskan arti filosofis peluang; ia menetapkan aritmetika yang
harus dipatuhi setiap perhitungan.

\kamus{peluang}{probability}{bobot pada setiap kemungkinan, dengan total bobot seluruh
kemungkinan sama dengan satu}

\index{peluang}
\index{ruang sampel}

\begin{definisi}[Ruang peluang]
Misalkan $\Omega$ adalah himpunan hasil, $\mathcal{F}$ kumpulan kejadian, yaitu
himpunan bagian $\Omega$ yang memuat $\Omega$ dan tertutup terhadap komplemen serta
gabungan terbilang, dan $P\colon\mathcal{F}\to\mathbb{R}$. Fungsi $P$ disebut
\emph{peluang} bila
\begin{enumerate}
  \item $P(A)\ge 0$ untuk setiap $A\in\mathcal{F}$;
  \item $P(\Omega)=1$;
  \item bila $A_{1},A_{2},\dots$ saling lepas, maka
        $P\bigl(\bigcup_{i\ge 1}A_{i}\bigr)=\sum_{i\ge 1}P(A_{i})$.
\end{enumerate}
\end{definisi}

Aksioma ketiga, aditif terbilang, membedakan teori peluang modern dari perhitungan
untung-untungan abad ketujuh belas. Tanpanya, deret peluang seperti
$\sum_{n\ge 1}2^{-n}=1$ tidak punya dasar, padahal deret itulah yang muncul ketika kita
bertanya pada lemparan ke berapa keberhasilan pertama terjadi. Aksioma ini juga yang
membuat peluang dapat dipakai pada limit, persoalan Bab 4 dan Bab 5.

Tiga akibat langsung berikut dipakai terus-menerus, jadi layak dibuktikan sekali.

\begin{proposisi}[Akibat aksioma]
\index{aksioma peluang}
Untuk kejadian $A$ dan $B$ dalam ruang peluang berlaku
$P(A^{c})=1-P(A)$, $P(\varnothing)=0$, dan bila $A\subseteq B$ maka
$P(A)\le P(B)$. Selain itu
$P(A\cup B)=P(A)+P(B)-P(A\cap B)$.
\end{proposisi}

\begin{proof}[Bukti]
Karena $A$ dan $A^{c}$ saling lepas dan gabungannya $\Omega$, aksioma kedua dan ketiga
memberi $1=P(\Omega)=P(A)+P(A^{c})$. Untuk yang kedua, ambil $A=\varnothing$ pada hasil
pertama. Untuk yang ketiga, tulis $B$ sebagai gabungan saling lepas $A$ dan
$B\setminus A$; maka $P(B)=P(A)+P(B\setminus A)\ge P(A)$ karena peluang tidak negatif.
Untuk yang terakhir, pecah $A\cup B$ menjadi tiga bagian saling lepas $A\setminus B$,
$A\cap B$, dan $B\setminus A$. Dengan cara yang sama $A$ dan $B$ masing-masing terpecah
menjadi dua bagian, sehingga
\[
P(A\cup B)=P(A\setminus B)+P(A\cap B)+P(B\setminus A)
=P(A)+P(B)-P(A\cap B).\qedhere
\]
\end{proof}

Sifat terakhir itu, inklusi--eksklusi, adalah alasan mengapa ``atau'' tidak boleh
dijumlahkan begitu saja. Contohnya, di satu kelas $30$ mahasiswa, $18$ menyukai kopi
dan $12$ menyukai teh, serta $8$ menyukai keduanya. Maka peluang seorang mahasiswa yang
dipilih acak menyukai kopi atau teh adalah
\[
\frac{18+12-8}{30}=\frac{22}{30}=0.7333,
\]
bukan $30/30=1$ seperti yang diberikan penjumlahan naif $18+12=30$. Delapan orang
dihitung dua kali, dan inklusi--eksklusi menguranginya sekali.

Untuk tiga kejadian, rumusnya memperpanjang dengan tanda berganti:
\[
P(A\cup B\cup C)=P(A)+P(B)+P(C)-P(A\cap B)-P(A\cap C)-P(B\cap C)+P(A\cap B\cap C).
\]
Buktinya menerapkan versi dua kejadian dua kali, dan itu latihan pertama bab ini. Pola
tandanya bukan hiasan: bagian yang terlalu banyak dikurangi harus ditambahkan kembali.

Satu peringatan untuk Bab 18. Pada ruang berhingga yang setiap hasilnya berpeluang
positif, $P(A)=0$ berarti $A$ kosong; pada ruang kontinu tidak, sebab satu titik dapat
berpeluang nol tanpa mustahil. Bab ini bekerja di ruang berhingga.

\begin{contoh}
Dua dadu dilempar. Berapa peluang jumlah matanya $8$ atau kedua matanya sama? Dengan
$36$ hasil yang setara, kejadian pertama memuat $5$ hasil, kejadian kedua memuat $6$,
dan irisannya satu hasil, yaitu $(4,4)$. Jadi peluangnya
$(5+6-1)/36=10/36=0.2778$. Tanpa mengurangkan irisan, jawabannya $11/36=0.3056$, terlalu
besar sebesar $1/36$.
\end{contoh}

\section{Pencacahan dan peluang yang setara}
\tingkat{menengah}

Bila semua hasil sama mungkin, peluang sebuah kejadian tinggal nisbah banyak hasil yang
menguntungkan terhadap banyak hasil seluruhnya. Kalimat sehari-harinya: membagi jumlah
cara yang kita inginkan dengan jumlah cara yang ada. Definisi ini bukan hukum alam; ia
sebuah asumsi tentang simetri, dan asumsi itu harus diperiksa, bukan diandaikan.

\index{pencacahan}
\index{kombinasi}

\kamus{kombinasi}{combination}{pilihan beberapa benda tanpa memperhatikan urutannya,
seperti mengambil lima kartu dari satu tumpukan}

\begin{proposisi}[Aturan perkalian dan pencacahan]
\index{aturan perkalian}
Bila $|A|=m$ dan $|B|=n$, maka $|A\times B|=mn$. Banyak barisan panjang $k$ dari $n$
benda berbeda tanpa pengulangan adalah $n!/(n-k)!$, dan banyak himpunan bagian
berukuran $k$ adalah
\[
\binom{n}{k}=\frac{n!}{k!\,(n-k)!}.
\]
Khususnya $\sum_{k=0}^{n}\binom{n}{k}=2^{n}$.
\end{proposisi}

\begin{proof}[Bukti]
Untuk hasilkali, setiap pasangan $(a,b)$ ditentukan oleh satu pilihan dari $m$ dan satu
pilihan dari $n$. Untuk barisan, ada $n$ pilihan untuk tempat pertama, $n-1$ untuk
tempat kedua, dan seterusnya sampai $n-k+1$. Untuk himpunan bagian, setiap himpunan
berukuran $k$ dapat diurutkan dengan $k!$ cara, sehingga banyak barisan dibagi $k!$
memberi banyak himpunan. Identitas terakhir menghitung semua himpunan bagian dengan dua
cara: menjumlahkan menurut ukuran, atau memilih masuk atau tidak untuk setiap benda.
\end{proof}

Pencacahan yang rapi sering menyingkat perhitungan panjang. Tangan penuh dalam poker,
tiga kartu seperingkat ditambah sepasang kartu seperingkat lain, memberi contoh bagus:
pilih peringkat trio ($13$ cara), pilih tiga dari empat jenisnya ($\binom{4}{3}=4$),
pilih peringkat pasangan ($12$ cara), pilih dua dari empat jenisnya ($\binom{4}{2}=6$).
Jumlahnya $13\cdot 4\cdot 12\cdot 6=3{,}744$ dari $\binom{52}{5}=2{,}598{,}960$ tangan,
sehingga peluangnya $0.001441$, yaitu satu dari $694.2$ tangan. Dalam $694$ tangan,
banyaknya tangan penuh yang diharapkan $0.9998$, jadi kita menunggu sekitar tujuh ratus
tangan untuk melihat yang pertama.

Pencacahan juga menjelaskan kejadian yang terasa ajaib. Ambil masalah ulang tahun:
$k$ orang di satu ruangan, setiap hari lahir setara di antara $365$ hari, dan hari
lahir antarorang saling bebas. Peluang tidak ada dua orang berbagi hari lahir adalah
\[
P(\text{tidak ada pasangan})=\frac{365\cdot 364\cdots (365-k+1)}{365^{k}},
\]
sehingga peluang ada pasangan adalah pelengkapnya.

\index{paradoks ulang tahun}

\begin{center}
\small
\begin{tabular}{rrr}
\hline
$k$ & $P(\text{tidak ada pasangan})$ & $P(\text{ada pasangan})$ \\
\hline
$10$ & $0.8831$ & $0.1169$ \\
$23$ & $0.4927$ & $0.5073$ \\
$41$ & $0.0968$ & $0.9032$ \\
$50$ & $0.0296$ & $0.9704$ \\
$57$ & $0.0099$ & $0.9901$ \\
\hline
\end{tabular}
\end{center}

Pada $k=23$ peluangnya melewati setengah, dan pada $k=50$ sudah $0.9704$. Sumber
kebingungannya satuan yang salah: yang bertabrakan bukan orang, melainkan pasangan orang,
dan $23$ orang membentuk $\binom{23}{2}=253$ pasangan. Setiap pasangan cocok dengan
peluang $1/365$, sehingga banyak pasangan yang diharapkan $253/365=0.69315$. Bila
tabrakan dianggap jarang dan saling bebas, peluang tidak ada tabrakan mendekati
$e^{-0.69315}=0.5000$, berselisih $0.0073$ dari nilai eksak $0.4927$. Pendekatan itu
bukan bukti; ia menjelaskan mengapa ambangnya jatuh di sekitar dua puluh tiga.

Pencacahan di atas mengandaikan setiap hari lahir setara. Pada populasi nyata
ketidaksetaraan itu justru menaikkan peluang tabrakan; asumsi simetri di sini adalah
pendekatan yang perlu diperiksa.

Pencacahan tidak dapat menjawab pertanyaan klinik tadi. Hasil ``sakit'' dan ``sehat''
tidak sama mungkin: nisbahnya satu berbanding sembilan ratus sembilan puluh sembilan.
Yang dibutuhkan bukan cara menghitung hasil, melainkan cara memakai informasi bahwa
seseorang sudah dinyatakan positif. Itulah peluang bersyarat.

\section{Peluang bersyarat dan aturan perkalian}
\tingkat{menengah}

Informasi baru mempersempit dunia. Kalau kita tahu $B$ terjadi, dunia yang relevan bukan
lagi $\Omega$ melainkan $B$, dan semua bobot harus dibaca ulang di dalam dunia yang lebih
sempit itu. Kalimat sehari-harinya: peluang bersyarat adalah peluang yang dihitung ulang
setelah sebagian kemungkinan dicoret.

\kamus{peluang bersyarat}{conditional probability}{peluang yang dihitung ulang setelah
kita tahu sebagian kemungkinan sudah tidak berlaku}

\index{peluang bersyarat}

\begin{definisi}[Peluang bersyarat]
Bila $A$ dan $B$ kejadian dengan $P(B)>0$, peluang $A$ dengan syarat $B$ adalah
\[
P(A\mid B)=\frac{P(A\cap B)}{P(B)}.
\]
\end{definisi}

Pembagian dengan $P(B)$ adalah pilihan yang membuat peluang bersyarat berjumlah satu di
atas $B$. Bila kita juga menuntut $P(B\mid B)=1$ dan nisbah antar-kejadian di dalam $B$
tidak berubah, pembagian itulah satu-satunya pilihan; karena itu definisi ini bukan
konvensi sembarangan.

Penyempitan itu dapat dilihat langsung pada tabel silang. Gambar~\ref{fig:silang}
menyusun populasi klinik menjadi empat sel. Sebelum ada hasil, penyebutnya
seluruh $100{,}000$ orang dan $P(S)=100/100{,}000=0.001$. Setelah hasil positif
diketahui, yang masih dihitung hanya kolom $+$ yang berisi $1{,}098$ orang, dan
$P(S\mid+)=99/1{,}098=0.0902$. Orangnya tidak berubah; yang berubah hanya himpunan yang
masih boleh dihitung.

    \begin{figure}[htbp]
\centering
\begin{tikzpicture}[font=\small, x=1cm, y=1cm]
  \fill[accent!10] (1.95,1.05) rectangle (3.25,-1.1);
  \draw[layergrey, thick] (1.9,1.05)--(7.3,1.05);
  \draw[layergrey] (1.9,0.35)--(7.3,0.35);
  \draw[layergrey] (1.9,-0.35)--(7.3,-0.35);
  \draw[layergrey, thick] (1.9,-1.1)--(7.3,-1.1);
  \node[anchor=east] at (1.75,1.4) {hasil uji};
  \node at (2.6,1.4) {$+$};
  \node at (4.3,1.4) {$-$};
  \node at (6.5,1.4) {jumlah};
  \node[anchor=east] at (1.75,0.7) {sakit $S$};
  \node at (2.6,0.7) {99};
  \node at (4.3,0.7) {1};
  \node at (6.5,0.7) {100};
  \node[anchor=east] at (1.75,0) {sehat $S^{c}$};
  \node at (2.6,0) {999};
  \node at (4.3,0) {98\,901};
  \node at (6.5,0) {99\,900};
  \node[anchor=east] at (1.75,-0.7) {jumlah};
  \node at (2.6,-0.7) {1\,098};
  \node at (4.3,-0.7) {98\,902};
  \node at (6.5,-0.7) {100\,000};
  \draw[->, accent, thick] (2.6,-1.65)--(2.6,-1.15);
  \node[anchor=north, align=center, text width=10.4cm, font=\footnotesize] at (4.6,-1.95)
    {setelah $+$ diketahui, hanya kolom ini yang dihitung:
     $99/1{,}098=11/122=0.0902$};
\end{tikzpicture}
\caption{Tabel silang uji saring pada $100{,}000$ orang. Baris adalah keadaan
sebenarnya, kolom adalah hasil uji. Setelah hasil positif diketahui, ruang sampel
menyempit dari seluruh tabel ke kolom $+$ yang berisi $1{,}098$ orang, dan $99$
di antaranya benar-benar sakit.}
\label{fig:silang}
\end{figure}    

\begin{proposisi}[Peluang bersyarat adalah peluang]
\index{peluang bersyarat}
Untuk $B$ tetap dengan $P(B)>0$, fungsi $A\mapsto P(A\mid B)$ memenuhi ketiga aksioma
peluang pada $\mathcal{F}$.
\end{proposisi}

\begin{proof}[Bukti]
Ketidaknegatifan berasal dari $P(A\cap B)\ge 0$. Norma berasal dari $P(\Omega\mid B)=1$,
sebab $P(B)/P(B)=1$. Untuk aditivitas, bila $A_{1},A_{2},\dots$ saling lepas
maka $A_{i}\cap B$ juga saling lepas, sehingga
\[
P\Bigl(\bigcup_{i}A_{i}\mid B\Bigr)
=\frac{P\bigl(\bigcup_{i}(A_{i}\cap B)\bigr)}{P(B)}
=\frac{\sum_{i}P(A_{i}\cap B)}{P(B)}
=\sum_{i}P(A_{i}\mid B).\qedhere
\]
\end{proof}

Karena itu semua akibat aksioma berlaku untuk peluang bersyarat: komplemen, monotonisitas,
dan inklusi--eksklusi tetap sah bila setiap $P$ diganti $P(\cdot\mid B)$. Yang tidak
berlaku adalah menukar arah syaratnya.

Dari definisi langsung lahir aturan perkalian
\[
P(A\cap B)=P(A\mid B)\,P(B)=P(B\mid A)\,P(A),
\]
yang memperpanjang menjadi
$P(A_{1}\cap\cdots\cap A_{n})=P(A_{1})P(A_{2}\mid A_{1})\cdots
P(A_{n}\mid A_{1}\cap\cdots\cap A_{n-1})$. Aturan ini membaca percobaan bertahap
sebagai perkalian peluang di setiap tahap.

Satu tumpukan kartu $52$ memberi contoh arah yang mudah tertukar. Misalkan $A$ kejadian
kartu pertama hati dan $B$ kejadian kartu kedua merah. Maka
\[
P(A\cap B)=\frac{13\cdot 25}{52\cdot 51}=\frac{25}{204},
\]
sehingga
\[
P(A\mid B)=\frac{25/204}{1/2}=\frac{25}{102}=0.2451,
\qquad
P(B\mid A)=\frac{25/204}{1/4}=\frac{25}{51}=0.4902 .
\]
Dua angka yang berbeda dua kali lipat, dari data yang sama. Kesalahannya bukan pada
aritmetika, melainkan pada penempatan kejadian di kiri dan kanan garis tegak.

\begin{bacakeras}
Dua suara. Yang pertama belum tahu, yang kedua sudah tahu. Bacakan bergantian.

Suara pertama: ``Duniaku seluas lapangan. Semua pintu mungkin. Aku menebak dengan mata
tertutup.''

Suara kedua: ``Duniaku satu kamar. Satu pintu sudah terbuka, dan aku hanya menghitung
yang ada di dalamnya.''

Suara pertama: ``Kalau aku tahu sesuatu, apakah peluangku berubah?''

Suara kedua: ``Bukan peluangmu yang berubah. Dunia yang kamu hitung yang berubah. Dunia
mengecil setelah kita tahu sesuatu: bobot yang tadinya tersebar sekarang berkumpul di
tempat yang masih mungkin. Penyebutnya menyempit, dan angka yang keluar terdengar baru
walaupun aritmetikanya tidak pernah berubah sejak awal.''

Suara pertama: ``Jadi yang bergerak penyebutnya?''

Suara kedua: ``Yang bergerak penyebutnya, dan yang tinggal dihitung ulang adalah sisa
yang masih boleh ikut.''

Ulangi bersama, tiga kali, makin pelan: ``Dunia mengecil, bukan peluangnya yang
berkhianat. Yang berubah hanya yang masih boleh dihitung.''
\end{bacakeras}

Alat kedua memecah populasi menjadi beberapa kelompok yang tidak bertumpuk. Bila
$B_{1},\dots,B_{k}$ saling lepas dan gabungannya $\Omega$, maka
\[
P(A)=\sum_{i=1}^{k}P(A\mid B_{i})\,P(B_{i}).
\]

\begin{proposisi}[Hukum peluang total]
\index{peluang total}
Bila $B_{1},\dots,B_{k}$ saling lepas, gabungannya $\Omega$, dan setiap $P(B_{i})>0$,
maka $P(A)=\sum_{i}P(A\mid B_{i})P(B_{i})$ untuk setiap kejadian $A$.
\end{proposisi}

\begin{proof}[Bukti]
Kejadian $A\cap B_{1},\dots,A\cap B_{k}$ saling lepas dan gabungannya $A$. Jadi
$P(A)=\sum_{i}P(A\cap B_{i})$, dan setiap suku ditulis ulang dengan aturan perkalian
sebagai $P(A\mid B_{i})P(B_{i})$.
\end{proof}

Hukum ini tampak seperti formalitas sampai ia dipakai pada contoh klinik: ia yang
mengubah dua angka iklan menjadi satu angka yang berguna, dan itulah isi bagian
berikutnya.

Sebelum itu, satu contoh bahwa kalimat Indonesia dan kejadian matematis tidak selalu
sepadan. Sebuah keluarga punya dua anak, setiap anak laki-laki atau perempuan dengan
peluang sama, dan jenis kelamin keduanya saling bebas. Bila kita tahu \emph{paling sedikit
satu} anak laki-laki, peluang keduanya laki-laki $1/3=0.3333$: dari empat susunan setara
tersisa tiga, dan hanya satu yang cocok. Bila yang kita tahu \emph{anak yang lebih tua}
laki-laki, peluangnya $1/2=0.5$. Dua kalimat yang mirip memberi jawaban berbeda, sebab
yang pertama mengecualikan satu susunan dan yang kedua dua.

\section{Aturan Bayes dan tafsir hasil uji}
\tingkat{menengah}

Aturan Bayes membalik arah syarat: dari peluang hasil bila sakit, ia menghitung peluang
sakit bila hasil. Kalimat sehari-harinya: dari semua cara sebuah hasil dapat muncul,
berapa bagian yang lewat jalan yang kita tanyakan.

\index{aturan bayes}
\index{kelaziman}

\begin{teorema}[Aturan Bayes]
Bila $B_{1},\dots,B_{k}$ saling lepas dengan gabungan $\Omega$ dan $P(B_{i})>0$, maka
untuk setiap kejadian $A$ dengan $P(A)>0$,
\[
P(B_{j}\mid A)=\frac{P(A\mid B_{j})\,P(B_{j})}
{\sum_{i=1}^{k}P(A\mid B_{i})\,P(B_{i})}.
\]
\end{teorema}

\begin{proof}[Bukti]
Menurut definisi, $P(B_{j}\mid A)=P(A\cap B_{j})/P(A)$. Pembilangnya adalah
$P(A\mid B_{j})P(B_{j})$ menurut aturan perkalian, dan penyebutnya $P(A)$ menurut hukum
peluang total. Jadi rumus itu tidak memuat asumsi baru; ia definisi peluang bersyarat
yang ditulis dua kali dengan arah berbeda.
\end{proof}

Karena itu tidak ada ``metode Bayes'' yang berdiri di atas aksioma; yang ada adalah
pembukuan. Pembilang adalah jalur yang kita tanyakan, penyebut adalah jumlah semua jalur
yang menghasilkan hasil yang sama.

\begin{resep}
\begin{enumerate}
  \item Namai dulu dua hal: hipotesis $H$ yang ditanyakan dan bukti $E$ yang sudah
        terlihat. Tulis pertanyaannya sebagai $P(H\mid E)$, bukan $P(E\mid H)$.
  \item Tulis pembilang sebagai satu jalur saja: $P(E\mid H)\,P(H)$. Itu bobot dunia
        tempat $H$ benar.
  \item Tulis penyebut sebagai jumlah semua jalur yang menghasilkan $E$, termasuk jalur
        lewat $H^{c}$: $P(E\mid H)P(H)+P(E\mid H^{c})P(H^{c})$.
  \item Bagi, lalu periksa akal: hasilnya antara $0$ dan $1$, dan lebih kecil daripada
        $P(E\mid H)$ bila $H$ jarang.
\end{enumerate}
Peringatan urutan: pembilang dan penyebut sama-sama memakai $P(E\mid\cdot)$. Menulis
$P(\cdot\mid E)$ di penyebut adalah kesalahan yang paling sering terjadi, dan besarnya
angka tidak memberi tahu bahwa arahnya sudah tertukar.
\end{resep}

Pohon pada Gambar~\ref{fig:pohon} menuliskan seluruh jalan menuju hasil positif. Setiap
daun adalah darab peluang di cabangnya, dan dua daun bertanda $+$ berjumlah
$P(+)=0.01098$. Pembilang Bayes hanyalah daun pertama, sedangkan penyebutnya jumlah
kedua daun itu.

    \begin{figure}[htbp]
\centering
\begin{tikzpicture}[font=\small, every node/.style={inner sep=2pt}]
  \node (r) at (0,0) {awal};
  \node (s) at (2.2,1.7) {$S$};
  \node (sc) at (2.2,-1.7) {$S^{c}$};
  \node (sp) at (4.6,2.5) {$+$};
  \node (sm) at (4.6,0.9) {$-$};
  \node (scp) at (4.6,-0.9) {$+$};
  \node (scm) at (4.6,-2.5) {$-$};
  \draw[->] (r)--(s) node[midway, above left=-2pt] {$0.001$};
  \draw[->] (r)--(sc) node[midway, below left=-2pt] {$0.999$};
  \draw[->, accent, thick] (s)--(sp) node[midway, above left=-2pt] {$0.99$};
  \draw[->, layergrey] (s)--(sm) node[midway, below left=-2pt] {$0.01$};
  \draw[->, accent, thick] (sc)--(scp) node[midway, above left=-2pt] {$0.01$};
  \draw[->, layergrey] (sc)--(scm) node[midway, below left=-2pt] {$0.99$};
  \node[anchor=west, font=\footnotesize] at (5.0,2.5) {$0.001\cdot 0.99=0.00099$};
  \node[anchor=west, font=\footnotesize] at (5.0,0.9) {$0.001\cdot 0.01=0.00001$};
  \node[anchor=west, font=\footnotesize] at (5.0,-0.9) {$0.999\cdot 0.01=0.00999$};
  \node[anchor=west, font=\footnotesize] at (5.0,-2.5) {$0.999\cdot 0.99=0.98901$};
  \node[anchor=north, align=center, text width=10.6cm, font=\footnotesize] at (4.4,-3.0)
    {$P(+)=0.00099+0.00999=0.01098$, \quad $P(S\mid+)=0.00099/0.01098=0.0902$};
\end{tikzpicture}
\caption{Pohon dua tahap uji saring. Cabang pertama memisahkan sakit dari sehat,
cabang kedua memisahkan hasil positif dari negatif; darab di tiap daun adalah peluang
jalur itu, dan dua daun positif berjumlah $P(+)=0.01098$.}
\label{fig:pohon}
\end{figure}    

Kembali ke klinik. Tulis $S$ untuk kejadian sakit dan $+$ untuk hasil positif. Dari
iklan dan kelaziman penyakit kita punya $P(S)=0.001$, $P(+\mid S)=0.99$, dan
$P(+\mid S^{c})=0.01$. Hukum peluang total memberi
\[
P(+)=0.99\cdot 0.001+0.01\cdot 0.999=0.00099+0.00999=0.01098,
\]
sehingga
\[
P(S\mid +)=\frac{0.00099}{0.01098}=\frac{11}{122}=0.0902 .
\]
Suku positif palsu, $0.00999$, sepuluh kali lebih besar daripada suku deteksi benar,
$0.00099$; nisbahnya $999/99=10.09$. Bukan karena ujinya buruk: orang sehat $999$ kali
lebih banyak, $99{,}900$ berbanding $100$.

Perubahan kelaziman mengubah jawaban jauh lebih kuat daripada perubahan mutu uji.
Dengan sensitivitas dan spesifisitas tetap $0.99$, tabel berikut menunjukkan
$P(S\mid +)$ untuk beberapa kelaziman.

\begin{center}
\small
\begin{tabular}{rrrr}
\hline
$P(S)$ & $P(S\mid +)$ & $1/P(S\mid +)$ & tafsiran \\
\hline
$0.0001$ & $0.0098$ & $102.0$ & satu dari seratus dua \\
$0.001$ & $0.0902$ & $11.09$ & satu dari sebelas \\
$0.010$ & $0.5000$ & $2.00$ & satu dari dua \\
$0.100$ & $0.9167$ & $1.09$ & sebelas dari dua belas \\
\hline
\end{tabular}
\end{center}

Pada kelaziman $1$ persen, uji yang sama memberi $P(S\mid+)=1/2$ tepat, bukan $0.99$.
Pada kelaziman $10$ persen, $P(S\mid+)=11/12=0.9167$. Angka-angka ini menjelaskan
kebijakan kesehatan yang sering dikritik: uji saring massal untuk penyakit langka
menghasilkan banyak positif palsu, dan karena itu hasil positif biasanya diikuti uji
kedua, bukan pengobatan.

Uji kedua memakai aturan yang sama, dengan prior yang sudah diperbarui. Bila uji kedua
bebas dari yang pertama dan bermutu sama, maka
\[
P(S\mid +,+)=\frac{0.99\cdot (11/122)}{0.99\cdot (11/122)+0.01\cdot (111/122)}
=\frac{10.89}{10.89+1.11}=\frac{363}{400}=0.9075 .
\]
Perhatikan bahwa $0.9075$ bukan $0.99$ dan bukan $0.9999$. Dua uji positif menaikkan
peluang dari $0.0902$ ke $0.9075$, hampir sepuluh kali, tetapi belum memberi kepastian.
Yang berubah hanya kelaziman di dalam kelompok yang sudah positif; mutu uji tetap. Dengan
mutu $0.99$ ini, posterior melampaui $0.99$ hanya bila kelaziman awal sudah melewati
$1/2$.

Satu hal teknis: $P(S\mid+)=0.0902$ berlaku untuk populasi berkelaziman $1/1000$, bukan
untuk pasien yang datang dengan gejala. Gejala mengubah kelaziman lebih dahulu, dan
itulah sebabnya dokter menanyakan riwayat sebelum menafsirkan hasil laboratorium.

Contoh terakhir bagian ini melepaskan Bayes dari uji klinik. Tiga kantong berisi sepuluh
kelereng: kantong $A$ satu merah, $B$ lima merah, $C$ sembilan merah. Kita memilih
kantong secara acak, mengambil satu kelereng, dan yang keluar merah. Sebelum pengambilan
ketiga kantong berpeluang $1/3$; sesudahnya $0.0667$, $0.3333$, dan $0.6$, seperti
Gambar~\ref{fig:bayes-batang}. Tidak ada bobot yang hilang; bobot hanya berpindah ke
kantong yang lebih mudah memberi bukti itu.

    \begin{figure}[htbp]
\centering
\begin{tikzpicture}[font=\small, x=1cm, y=1cm]
  \draw[->, layergrey] (0,0)--(7.4,0);
  \draw[->, layergrey] (0,0)--(0,4.3);
  \foreach \y/\lab in {0/0, 1.2/0.2, 2.4/0.4, 3.6/0.6} {
    \draw[layergrey] (0,\y)--(-0.08,\y);
    \node[anchor=east, font=\footnotesize] at (-0.12,\y) {\lab};
  }
  \fill[layerindigo] (0.65,0) rectangle (1.20,2.0);
  \fill[layerteal] (1.40,0) rectangle (1.95,2.0);
  \fill[accent] (2.15,0) rectangle (2.70,2.0);
  \fill[layerindigo] (4.35,0) rectangle (4.90,0.4);
  \fill[layerteal] (5.10,0) rectangle (5.65,2.0);
  \fill[accent] (5.85,0) rectangle (6.40,3.6);
  \draw[dashed, layergrey] (0,2.0)--(6.9,2.0);
  \node[font=\footnotesize, anchor=south] at (0.92,2.02) {$1/3$};
  \node[font=\footnotesize, anchor=south] at (1.67,2.02) {$1/3$};
  \node[font=\footnotesize, anchor=south] at (2.42,2.02) {$1/3$};
  \node[font=\footnotesize, anchor=south] at (4.62,0.44) {$0.0667$};
  \node[font=\footnotesize, anchor=south] at (5.37,2.02) {$0.3333$};
  \node[font=\footnotesize, anchor=south] at (6.12,3.62) {$0.6$};
  \node[font=\footnotesize, anchor=north] at (0.92,-0.08) {$A$};
  \node[font=\footnotesize, anchor=north] at (1.67,-0.08) {$B$};
  \node[font=\footnotesize, anchor=north] at (2.42,-0.08) {$C$};
  \node[font=\footnotesize, anchor=north] at (4.62,-0.08) {$A$};
  \node[font=\footnotesize, anchor=north] at (5.37,-0.08) {$B$};
  \node[font=\footnotesize, anchor=north] at (6.12,-0.08) {$C$};
  \node[font=\footnotesize] at (1.67,-0.85) {sebelum};
  \node[font=\footnotesize] at (5.37,-0.85) {sesudah};
\end{tikzpicture}
\caption{Pembaruan Bayes pada tiga kantong. Sebelum ada bukti ketiga kantong berbagi
bobot $1/3$; setelah satu kelereng merah terlihat, bobot berpindah menjadi $0.0667$,
$0.3333$, dan $0.6$. Kantong yang paling mudah menghasilkan bukti itu naik paling
tinggi.}
\label{fig:bayes-batang}
\end{figure}    

\section{Independensi, dan mengapa saling lepas bukan independen}
\tingkat{lanjut}

Dua kejadian disebut saling lepas bila keduanya tidak dapat terjadi bersamaan. Dua
kejadian disebut independen bila mengetahui yang satu tidak mengubah peluang yang lain.
Kedua kalimat itu kedengarannya serupa, dan menukarnya adalah kekeliruan paling mahal
di bab ini. Kalimat sehari-harinya: saling lepas berarti tidak pernah bersama;
independen berarti tidak saling memberi kabar.

\kamus{independen}{independent}{kejadian yang tidak saling memberi kabar: mengetahui
satu tidak mengubah peluang yang lain}

\begin{untukanak}
Ada dua hal yang sering tertukar, dan bedanya dapat dilihat tanpa angka. Di satu
sekolah, dua siswa tidak pernah masuk kelas bersamaan: yang satu belajar pagi, yang satu
belajar siang. Kalau kita tahu Andi sudah masuk, kita tahu Budi belum masuk. Mereka tidak
pernah bersama, dan justru karena itu kabar tentang yang satu memberi kabar tentang yang
lain.

Sekarang dua siswa lain. Keduanya masuk kelas yang sama, tetapi duduk berjauhan dan
mengerjakan ujian yang berbeda. Nilai mereka tidak saling memberi kabar: tahu nilai Sari
tidak menolong kita menebak nilai Tono.

Kelompok pertama disebut tidak pernah bersama. Kelompok kedua disebut tidak saling
memberi kabar. Keduanya kedengaran seperti ``tidak ada hubungan'', padahal yang pertama
justru hubungan yang paling kuat: kabar tentang yang satu menentukan yang lain. Di bab
ini kita akan tahu mengapa dua hal yang tidak pernah bersama tidak mungkin saling memberi
kabar, dan mengapa menukar keduanya membuat hitungan meleset bukan sedikit.
\end{untukanak}

\index{independensi}
\index{kejadian saling lepas}

\begin{definisi}[Independensi dua kejadian]
Kejadian $A$ dan $B$ \emph{independen} bila
\[
P(A\cap B)=P(A)\,P(B).
\]
\end{definisi}

Definisi ini simetris, dan bila $P(B)>0$ ia setara dengan $P(A\mid B)=P(A)$. Perhatikan
bahwa definisi itu tidak menyebut sebab akibat: dua kejadian dapat independen tanpa satu
menyebabkan yang lain, dan dapat bergantung tanpa satu menyebabkan yang lain.

\begin{proposisi}[Saling lepas bukan independen]
\index{kejadian saling lepas}
Bila $A$ dan $B$ saling lepas dan $P(A)>0$ serta $P(B)>0$, maka $A$ dan $B$ tidak
independen. Lebih tepatnya $P(A\cap B)=0<P(A)P(B)$, dan $P(A\mid B)=0<P(A)$.
\end{proposisi}

\begin{proof}[Bukti]
Saling lepas berarti $A\cap B=\varnothing$, sehingga $P(A\cap B)=P(\varnothing)=0$.
Karena $P(A)P(B)>0$, kesamaan definisi independensi gagal. Untuk pernyataan bersyarat,
$P(A\mid B)=0/P(B)=0$, yang lebih kecil daripada $P(A)>0$.
\end{proof}

Kekeliruannya bergerak dua arah, dan kedua arah itu dapat dihitung.

Ambil satu dadu dan dua pasang kejadian dengan marginal sama tetapi hubungan berbeda.
Pasangan pertama: $A=\{1,2\}$ dan $B=\{3,4\}$, saling lepas, dengan $P(A)=P(B)=1/3$.
Pasangan kedua: $C=\{1,2,3\}$ dan $D=\{1,4\}$, dengan $P(C)=1/2$ dan $P(D)=1/3$. Untuk
pasangan pertama,
\[
P(A)P(B)=\frac{1}{9}=0.1111,\qquad P(A\cap B)=0,\qquad P(A\mid B)=0,
\]
jadi selisih definisi independensi $0.1111$ dan selisih peluang bersyaratnya $0.3333$,
sebab $P(A)=1/3$. Untuk pasangan kedua,
\[
P(C)P(D)=\frac{1}{6}=0.1667=P(C\cap D),\qquad P(C\mid D)=\frac{1}{2}=P(C),
\]
jadi kedua selisihnya nol.

Yang penting bukan angka $1/6$ itu, melainkan bahwa dua pasangan dapat bermarginal sama
dengan hubungan berlawanan total. Pasangan $E=\{1,2,3\}$ dan $F=\{4,6\}$ juga bermarginal
$1/2$ dan $1/3$ tetapi saling lepas. Marginal tidak menentukan hubungan; hanya irisan
yang menentukannya, sehingga setiap klaim independensi adalah klaim tambahan tentang
dunia.

Untuk tiga kejadian, independensi berpasangan tidak cukup. Kejadian $A,B,C$ disebut
saling independen bila $P(A\cap B)=P(A)P(B)$, $P(A\cap C)=P(A)P(C)$,
$P(B\cap C)=P(B)P(C)$, dan $P(A\cap B\cap C)=P(A)P(B)P(C)$.

\begin{contoh}[Berpasangan tetapi tidak saling independen]
\index{independensi berpasangan}
Dua koin dilempar. Misalkan $A$ koin pertama gambar, $B$ koin kedua gambar, dan $C$
kedua koin sama; masing-masing berpeluang $1/2$. Setiap
pasangan beririsan satu hasil, jadi $P(A\cap B)=P(A\cap C)=P(B\cap C)=1/4$, yang sama
dengan $1/2\cdot 1/2$. Tetapi $A\cap B\cap C$ terjadi hanya bila keduanya gambar,
sehingga peluangnya $1/4$, sedangkan $P(A)P(B)P(C)=1/8$. Jadi ketiganya berpasangan
independen dan tidak saling independen; selisihnya $0.125$.
\end{contoh}

Peringatan terakhir: independensi tidak bertahan di bawah pengkondisian. Bila $A$ dan
$B$ independen berpeluang $1/2$, maka dengan syarat $A\cup B$, $P(A\mid A\cup B)=2/3$
sedangkan $P(A\cap B\mid A\cup B)=1/3\neq 4/9$. Menggabungkan dua kejadian independen
lalu mengkondisikan pada gabungannya justru menciptakan ketergantungan, sehingga kalimat
``$A$ dan $B$ independen'' selalu relatif terhadap ruang peluang yang dipakai.

\begin{tangan}
Bahan: dua belas kartu, enam pasang berperingkat sama (As sampai enam), setiap pasangan
satu merah dan satu hitam.

Langkah pertama, buat pasangan. Susun dua belas kartu terbuka dalam dua baris; letakkan
setiap pasangan seperingkat bersebelahan, merah di kiri, hitam di kanan. Periksa dengan
mata: peringkat tidak memberi kabar tentang warna, sebab setiap peringkat punya merah
dan hitam. Dua sifat itu tidak saling memberi kabar.

Langkah kedua, saring. Kocok, tutup semua, lalu buka satu kartu: misalnya empat hati.
Beri satu petunjuk, ``pasangan kartu ini berwarna hitam'', lalu sisihkan semua kartu
merah yang masih tertutup. Dari sebelas kartu tertutup, tinggal enam yang hitam.

Yang harus kaulihat: penyebutnya berubah dari sebelas menjadi enam, sedangkan kartunya
tidak berubah sama sekali. Peluang menebak pasangannya naik dari satu per sebelas menjadi
satu per enam. Dua sifat yang tadi tidak saling memberi kabar sekarang saling memberi
kabar, dan itu yang dilakukan satu petunjuk.
\end{tangan}

\section{Dua paradoks yang menuntut hitungan}
\tingkat{lanjut}

Paradoks dalam peluang hampir selalu tabrakan antara intuisi dan definisi, dan
tabrakan itu hanya selesai bila dihitung. Dua contoh berikut menunjukkan hal berbeda:
jawaban bergantung pada aturan main, dan pada satuan yang dihitung.

\index{monty hall}

Monty Hall. Tiga pintu tertutup; satu berisi mobil, dua berisi kambing. Kita memilih
satu pintu. Pembawa acara, yang tahu isi semua pintu, membuka satu pintu lain yang
berisi kambing, lalu menawarkan kesempatan berpindah ke pintu yang masih tertutup.
Apakah berpindah menaikkan peluang?

Hitung dengan hukum peluang total terhadap tiga kemungkinan letak mobil, masing-masing
berpeluang $1/3$, dan andaikan kita mula-mula memilih pintu $1$. Bila mobil di pintu
$1$, berpindah kalah. Bila mobil di pintu $2$, pembawa acara wajib membuka pintu $3$,
sehingga berpindah menang. Bila mobil di pintu $3$, ia membuka pintu $2$, dan berpindah
menang. Jadi berpindah menang pada dua dari tiga kemungkinan:
\[
P(\text{menang dengan berpindah})=\frac{2}{3},\qquad
P(\text{menang tanpa berpindah})=\frac{1}{3}.
\]
Selisihnya $1/3$, bukan nol.

Argumen ``dua pintu tersisa, jadi $1/2$'' keliru karena mengabaikan bahwa pembawa acara
tidak memilih pintu secara bebas. Ia dilarang membuka pintu berisi mobil dan dilarang
membuka pintu pilihan kita, sehingga tindakannya membawa informasi tentang pintu yang
tidak ia buka.

Informasi itu hilang bila aturannya berubah. Andaikan pembawa acara membuka satu dari
dua pintu lain secara acak, tanpa melihat isinya, dan permainan dibatalkan bila mobil
terlihat. Dengan syarat permainan berjalan sampai tawaran,
\[
P(\text{mobil di pintu }1\mid \text{kambing terlihat})
=\frac{1/3}{1/3+1/6+1/6}=\frac{1}{2},
\]
sebab bila mobil di pintu $2$ atau $3$, ia hanya punya peluang $1/2$ menampilkan kambing.
Pada aturan itu berpindah tidak menolong: ``peluang berpindah $2/3$'' bukan sifat pintu,
melainkan sifat protokol.

\index{paradoks ulang tahun}

Ulang tahun. Bagian sebelumnya sudah menghitung $0.5073$ untuk $23$ orang dan $0.9704$
untuk $50$ orang. Paradoksnya bukan pada angkanya, melainkan pada satuan: kita
membandingkannya dengan peluang satu orang tertentu cocok dengan kita, padahal yang
dihitung $253$ pasangan sekaligus. Kesalahan yang sama muncul pada tabel hash dan uji
kecocokan DNA.

Kekeliruan jaksa adalah kesalahan arah yang sama pada perkara pidana. Misalkan peluang
sebuah ciri muncul secara kebetulan satu per satu juta, dan ciri itu ditemukan pada
tersangka yang diambil dari basis data sejuta orang tanpa bukti lain. Maka
$P(\text{cocok}\mid\text{tidak bersalah})=10^{-6}$, tetapi
\[
P(\text{bersalah}\mid\text{cocok})
=\frac{10^{-6}\cdot 1}{10^{-6}\cdot 1+(1-10^{-6})\cdot 10^{-6}}=0.5000003,
\]
yaitu praktis setengah, bukan hampir pasti. Dua bilangan $10^{-6}$ itu besaran berbeda:
yang pertama peluang ciri secara kebetulan, yang kedua kelaziman tersangka di dalam basis
data. Menukarnya adalah wajah lain dari kekeliruan yang sama, kali ini di ruang
pengadilan.

\begin{tantangan}
Tiga kartu di atas meja. Kartu pertama merah di kedua sisinya, kartu kedua hitam di kedua
sisinya, kartu ketiga merah di satu sisi dan hitam di sisi lain. Kita mengocok ketiganya
tertutup, mengambil satu, dan meletakkannya di meja. Sisi yang terlihat merah. Berapa
peluang sisi di baliknya juga merah?

Hampir semua orang menjawab setengah: katanya tinggal dua kartu yang punya sisi merah,
jadi dua kemungkinan. Jawaban yang benar bukan setengah.

Petunjuk 1: jangan menghitung kartu, hitung sisi. Ada berapa sisi merah seluruhnya, dan
berapa di antaranya bersebelahan dengan sisi merah?

Petunjuk 2: sisi merah yang terlihat dapat berasal dari tiga sisi merah yang berbeda,
dan dua di antaranya berada pada kartu merah-merah.

Kalau mentok, kembali ke bagian peluang bersyarat dan tulis ulang ruang sampelnya sebagai
enam sisi, bukan tiga kartu.
\end{tantangan}

\section{Contoh hitung}
\tingkat{menengah}

Tiga hitungan berikut dikerjakan dengan pecahan eksak, supaya tidak ada pembulatan yang
menyembunyikan langkah. Semuanya sudah diperiksa dengan Python, dan angka yang muncul
adalah angka yang sama dengan yang dipakai di bagian sebelumnya.

Pertama, Bayes dua langkah. Dengan $P(S)=1/1000$, $P(+\mid S)=99/100$, dan
$P(+\mid S^{c})=1/100$, hukum peluang total memberi
\[
P(+)=\frac{99}{100}\cdot\frac{1}{1000}+\frac{1}{100}\cdot\frac{999}{1000}
=\frac{99+999}{100{,}000}=\frac{1098}{100{,}000}.
\]
Pembilangnya dua kelompok terpisah: $99$ orang sakit yang terdeteksi dan $999$ orang
sehat yang salah dituduh. Maka
\[
P(S\mid +)=\frac{99/100{,}000}{1098/100{,}000}=\frac{99}{1098}=\frac{11}{122}=0.090164 .
\]
Uji kedua bermutu sama memakai $11/122$ sebagai kelaziman baru dan memberi $0.9075$,
sama seperti hitungan di bagian sebelumnya. Penyebutnya $1200$, bukan angka berjuta-juta:
faktor $1/122$ muncul di kedua suku dan saling menghapus, tanda bahwa yang berubah hanya
kelaziman, bukan mutu uji.

Kedua, saling lepas versus independen, dengan marginal yang disamakan.

\begin{center}
\small
\begin{tabular}{lrrrr}
\hline
pasangan & $P(\text{pertama})$ & $P(\text{kedua})$ & $P(\text{irisan})$ & $P(\text{pertama}\mid\text{kedua})$ \\
\hline
$C=\{1,2,3\}$, $D=\{1,4\}$ & $0.5$ & $0.3333$ & $0.1667$ & $0.5$ \\
$E=\{1,2,3\}$, $F=\{4,6\}$ & $0.5$ & $0.3333$ & $0$ & $0$ \\
\hline
\end{tabular}
\end{center}

Pada baris pertama $P(C)P(D)=1/6=0.1667$ sama dengan irisan, jadi selisihnya $0$. Pada
baris kedua $P(E)P(F)=1/6$ sedangkan irisannya $0$, jadi selisihnya $0.1667$; peluang
bersyaratnya $0$ sedangkan $P(E)=0.5$. Siapa pun yang menyangka ``tidak bisa terjadi
bersama'' berarti ``tidak saling memengaruhi'' akan melaporkan $0.1667$ di tempat yang
seharusnya $0$.

Ketiga, Monty Hall dengan protokol yang dinyatakan. Bila pembawa acara selalu membuka
pintu kambing, peluang menang dengan berpindah $2/3=0.6667$ dan tanpa berpindah
$1/3=0.3333$. Bila ia membuka secara acak dan permainan dibatalkan ketika mobil
terlihat, permainan berjalan dengan peluang $2/3$, dan dengan syarat itu kedua strategi
bernilai $1/2=0.5$. Dua protokol, satu himpunan pintu, dua jawaban.

\kertas{Untuk uji dengan kelaziman $1/2000$, sensitivitas $0.95$, dan spesifisitas
$0.90$, hitung $P(S\mid+)$ dengan tabel $100{,}000$ orang seperti di atas, lalu hitung
lagi dengan spesifisitas $0.99$. Tulis kedua jawaban sebagai pecahan dan sebagai desimal
empat angka, dan jelaskan bagian mana yang berubah paling banyak: pembilang atau
penyebut. Setelah itu hitung $P(S\mid +,+)$ untuk uji kedua yang bermutu sama, dan
periksa apakah hasilnya masih di bawah $0.95$.}

\section{Lab komputasi: Monte Carlo sebagai alat verifikasi}
\tingkat{lanjut}

Semua angka di bab ini diturunkan, bukan diukur. Simulasi tidak menggantikan penurunan,
tetapi ia menangkap kesalahan yang paling sering terjadi: arah syarat yang tertukar,
penyebut yang keliru, dan asumsi independensi yang dipakai tanpa dasar. Bangkitkan
percobaan yang sama banyak kali, hitung nisbah keberhasilan, lalu bandingkan dengan nilai
eksak. Menurut hukum bilangan besar (Bab 19), nisbah itu menuju peluang sebenarnya; yang
tersisa hanya seberapa cepat, diukur oleh galat baku $\sqrt{p(1-p)/m}$ dengan $m$ banyak
percobaan yang menghasilkan kejadian itu.

\index{monte carlo}
\index{galat baku}

Benih tetap dipakai supaya keluarannya dapat diperiksa ulang kata demi kata. Benih tetap
tidak membuat hasilnya lebih benar; ia hanya membuatnya dapat diulang.

\begin{lab}
    \begin{verbatim}
import numpy as np

rng = np.random.default_rng(20261024)

print("== 1. uji diagnostik: prevalensi 0,1%, sens 99%, spek 99% ==")
n = 1_000_000
sakit = rng.random(n) < 0.001
positif = np.where(sakit, rng.random(n) < 0.99, rng.random(n) < 0.01)
m = int(positif.sum())
ppv = sakit[positif].mean()
galat = np.sqrt(ppv * (1 - ppv) / m)
eksak = 0.99 * 0.001 / (0.99 * 0.001 + 0.01 * 0.999)
print("sakit=%d  positif=%d  benar_positif=%d"
      % (sakit.sum(), m, int((sakit & positif).sum())))
print("PPV empiris=%.6f  eksak=%.6f  galat baku=%.6f  z=%.3f"
      % (ppv, eksak, galat, (ppv - eksak) / galat))

print("== 2. Monty Hall, 10^5 permainan ==")
N = 100_000
mobil = rng.integers(0, 3, N)
pilih = rng.integers(0, 3, N)
buka = np.where(pilih == mobil,
                (pilih + 1 + rng.integers(0, 2, N)) % 3,
                3 - pilih - mobil)
ganti = 3 - pilih - buka
menang_tetap = float((pilih == mobil).mean())
menang_ganti = float((ganti == mobil).mean())
se_ganti = np.sqrt(menang_ganti * (1 - menang_ganti) / N)
print("menang tanpa ganti=%.4f  menang ganti=%.4f  galat baku=%.5f"
      % (menang_tetap, menang_ganti, se_ganti))
print("eksak 1/3=%.6f  2/3=%.6f" % (1/3, 2/3))

print("== 3. pembawa acara membuka pintu secara acak ==")
buka_acak = (pilih + 1 + rng.integers(0, 2, N)) % 3
sah = buka_acak != mobil
ganti_acak = 3 - pilih - buka_acak
print("permainan sah=%d dari %d (%.4f)" % (sah.sum(), N, sah.mean()))
print("menang ganti (sah)=%.4f  tanpa ganti (sah)=%.4f"
      % (float((ganti_acak[sah] == mobil[sah]).mean()),
         float((pilih[sah] == mobil[sah]).mean())))

print("== 4. peluang ulang tahun, 20.000 kamar ==")
for k in (23, 50):
    hari = rng.integers(0, 365, size=(20000, k))
    hari.sort(axis=1)
    ada = (np.diff(hari, axis=1) == 0).any(axis=1)
    q = 1.0
    for i in range(k):
        q *= (365 - i) / 365
    f = float(ada.mean())
    se = np.sqrt(f * (1 - f) / 20000)
    print("k=%2d  empiris=%.4f  eksak=%.4f  galat baku=%.5f  z=%+.2f"
          % (k, f, 1 - q, se, (f - (1 - q)) / se))

print("== 5. saling lepas versus independen, 10^6 lemparan dadu ==")
N = 1_000_000
dadu = rng.integers(1, 7, N)
A = np.isin(dadu, [1, 2])
B = np.isin(dadu, [3, 4])
C = np.isin(dadu, [1, 2, 3])
D = np.isin(dadu, [1, 4])
print("lepas:  P(A)=%.4f P(B)=%.4f P(AnB)=%.4f P(A|B)=%.4f selisih=%.4f"
      % (A.mean(), B.mean(), (A & B).mean(), (A & B).sum() / B.sum(),
         A.mean() * B.mean() - (A & B).sum() / B.sum()))
print("bebas:  P(C)=%.4f P(D)=%.4f P(CnD)=%.4f P(C|D)=%.4f selisih=%.4f"
      % (C.mean(), D.mean(), (C & D).mean(), (C & D).sum() / D.sum(),
         abs(C.mean() - (C & D).sum() / D.sum())))
print("eksak lepas: 1/3, 1/3, 0, 0, 1/9=%.6f" % (1 / 9))
print("eksak bebas: 1/2, 1/3, 1/6, 1/2, 0")
\end{verbatim}    
\end{lab}

Keluaran sesungguhnya, apa adanya:

    \begin{verbatim}
== 1. uji diagnostik: prevalensi 0,1%, sens 99%, spek 99% ==
sakit=963  positif=11094  benar_positif=952
PPV empiris=0.085812  eksak=0.090164  galat baku=0.002659  z=-1.637
== 2. Monty Hall, 10^5 permainan ==
menang tanpa ganti=0.3325  menang ganti=0.6675  galat baku=0.00149
eksak 1/3=0.333333  2/3=0.666667
== 3. pembawa acara membuka pintu secara acak ==
permainan sah=66472 dari 100000 (0.6647)
menang ganti (sah)=0.4998  tanpa ganti (sah)=0.5002
== 4. peluang ulang tahun, 20.000 kamar ==
k=23  empiris=0.5019  eksak=0.5073  galat baku=0.00354  z=-1.54
k=50  empiris=0.9701  eksak=0.9704  galat baku=0.00120  z=-0.19
== 5. saling lepas versus independen, 10^6 lemparan dadu ==
lepas:  P(A)=0.3333 P(B)=0.3335 P(AnB)=0.0000 P(A|B)=0.0000 selisih=0.1112
bebas:  P(C)=0.5001 P(D)=0.3334 P(CnD)=0.1667 P(C|D)=0.5000 selisih=0.0001
eksak lepas: 1/3, 1/3, 0, 0, 1/9=0.111111
eksak bebas: 1/2, 1/3, 1/6, 1/2, 0
\end{verbatim}    

Empat hal perlu dibaca dari keluaran itu, dan dua di antaranya justru menunjukkan batas
simulasi.

Pertama, baris pertama bagian 1 menyebut $963$ orang sakit dari sejuta, sedangkan
harapannya $1{,}000$ dengan simpangan baku $\sqrt{10^{6}(0.001)(0.999)}=31.6$. Jadi
$963$ berada $1.17$ simpangan baku di bawah harapan: ayunan wajar, bukan kelalaian
program. Harapan hasil positif menjadi $963\cdot 0.99+999{,}037\cdot 0.01=10{,}943.7$,
dan yang teramati $11{,}094$ berjarak $1.4$ simpangan baku. Peluang bersyarat empiris
$0.0858$ berbeda dari eksak $0.0902$ sebesar $0.0044$, yaitu $1.64$ galat baku. Selisih
itu jauh lebih kecil daripada jarak $0.09$ ke $1$, sehingga tafsirnya tidak berubah.

Kedua, Monty Hall memberi $0.6675$ dan $0.3325$ dengan galat baku $0.0015$, sedangkan
selisih yang diharapkan antarstrategi $1/3$. Kesepakatan sampai tiga angka itu bukan
kebetulan; bentuk verifikasi yang paling berguna memang bukan memastikan angka ketiga,
melainkan memastikan selisih yang kita klaim ada dan bertanda benar.

Ketiga, bagian 3 dan 5 menunjukkan bahwa protokol dan hubungan antar-kejadian masuk ke
angka. Pembawa acara yang membuka pintu secara acak hanya membiarkan $66{,}472$ dari
$100{,}000$ permainan berjalan sampai tawaran, dan di antaranya berpindah menang
$0.4998$ sedangkan bertahan $0.5002$; selisih $0.0004$ itu seperempat galat baku.
Pasangan saling lepas memberi $P(A)P(B)-P(A\cap B)=0.1112$, cocok dengan $1/9$ sampai
tiga angka, sedangkan pasangan independen memberi $0.0001$, yaitu nol dalam ayunan acak.
Angka $0.1112$ adalah besar kesalahan orang yang menyangka dua kejadian saling lepas
pasti independen.

Yang tidak dibuktikan program: tidak satu pun dari kelima bagian itu membuktikan rumus.
Program hanya memeriksa rumus yang sudah dibuktikan, dengan sampel jauh lebih besar
daripada jangkauan tangan; Bab 19 menjelaskan mengapa itu sah.

\begin{untukdosen}
Urutan mengajar yang saya sarankan: mulai dari tabel $100{,}000$ orang, bukan dari
aksioma. Suruh mahasiswa mengisi empat sel: sakit dan positif, sakit dan negatif, sehat
dan positif, sehat dan negatif. Angka $99$, $1$, $999$, dan $98{,}901$ muncul sendiri,
dan dari situ $99/1{,}098$ menjadi kesimpulan kelas, bukan hasil ceramah. Sesudah itu
aksioma, akibatnya, pencacahan, peluang bersyarat, Bayes, lalu independensi. Monty Hall
paling baik ditaruh paling akhir, setelah mahasiswa mampu menuliskan protokol pembawa
acara sebagai kejadian bersyarat.

Dua miskonsepsi yang paling sering muncul. Pertama, ``saling lepas berarti independen''.
Gejalanya terlihat di aritmetika: mahasiswa menjumlahkan peluang dua kejadian yang
seharusnya dikalikan, atau sebaliknya, lalu tidak dapat menjelaskan mengapa hasilnya
melebihi satu. Kedua, ``sensitivitas sama dengan peluang sakit bila hasil positif''.
Gejalanya adalah jawaban $0.99$ untuk soal klinik, tanpa memakai kelaziman yang sudah
diberikan di soal.

Aktivitas papan tulis yang berhasil: dua mahasiswa memerankan pembawa acara dengan
protokol tertulis berbeda, satu selalu membuka pintu kambing dan satu membuka acak; kelas
mencatat dua puluh permainan lalu membandingkannya dengan $2/3$ dan $1/2$. Mahasiswa siap
lanjut bila dapat menyebut satu contoh dua kejadian independen dan satu contoh saling
lepas dengan marginal sama, serta menghitung peluang bersyarat dari tabel dua arah tanpa
membalik arah.
\end{untukdosen}

\section{Ringkasan}

\begin{ringkasan}
\begin{itemize}
  \item Peluang adalah fungsi pada kejadian yang tidak negatif, bernilai satu pada
        seluruh ruang, dan aditif terbilang pada kejadian saling lepas; komplemen,
        monotonisitas, dan inklusi--eksklusi diturunkan dari situ.
  \item Bila hasil sama mungkin, peluang adalah nisbah pencacahan, dengan aturan
        perkalian dan $\binom{n}{k}$; asumsi kesetaraan itu harus diperiksa.
  \item Peluang bersyarat $P(A\mid B)=P(A\cap B)/P(B)$ menghitung ulang di dalam $B$
        dan memenuhi ketiga aksioma; aturan perkalian dan hukum peluang total adalah
        akibatnya.
  \item Aturan Bayes bukan asumsi baru, melainkan definisi yang dibaca dua arah.
        Kelaziman menentukan jawaban lebih kuat daripada mutu uji: dengan sensitivitas
        dan spesifisitas $0.99$, kelaziman $0.1$ persen memberi $P(S\mid+)=0.0902$.
  \item Saling lepas dan independen berbeda; bila kedua peluang positif, kejadian
        saling lepas tidak mungkin independen, dan hanya irisan yang menentukan
        hubungan.
  \item Jawaban paradoks bergantung pada protokol dan satuan hitung: Monty Hall $2/3$
        pada protokol tetap dan $1/2$ pada protokol acak; ulang tahun $0.5073$ untuk
        $23$ orang karena yang dihitung $253$ pasangan. Simulasi berbenih tetap
        memeriksa, bukan membuktikan.
\end{itemize}
\end{ringkasan}

\section{Latihan}

\begin{latihan}
\begin{enumerate}
  \item Buktikan rumus inklusi--eksklusi untuk tiga kejadian dengan menerapkan versi
        dua kejadian pada $A\cup B$ dan $C$. Tunjukkan pada langkah mana suku
        $P(A\cap B\cap C)$ muncul, dan mengapa tandanya positif.
  \item Nilai klaim ini: ``Karena $P(A\mid B)=0$, $A$ dan $B$ saling lepas.'' Benar,
        salah, atau belum cukup bukti? Beri contoh dengan $P(A\mid B)=0$ tetapi
        $A\cap B\neq\varnothing$, atau jelaskan mengapa itu tidak mungkin.
  \item Sebuah uji berkelaziman $2$ persen, sensitivitas $0.95$, spesifisitas $0.90$.
        Hitung $P(S\mid+)$ dari tabel $10{,}000$ orang, lalu ulangi dengan
        spesifisitas $0.99$; nyatakan keduanya sebagai pecahan.
  \item Buktikan bahwa bila $A$ dan $B$ independen, maka $A$ dan $B^{c}$ juga
        independen. Gunakan $P(A)=P(A\cap B)+P(A\cap B^{c})$ dan jelaskan mengapa
        hasil yang sama berlaku untuk $A^{c}$ dan $B^{c}$.
  \item Hitung peluang ada pasangan hari lahir sama untuk $30$ orang, bandingkan
        dengan $1-e^{-\binom{30}{2}/365}$, lalu jelaskan sumber selisihnya.
  \item Buktikan hukum peluang total untuk $k=2$ langsung dari aksioma, tanpa memakai
        hukum untuk $k$ umum: tulis $A=(A\cap B)\cup(A\cap B^{c})$ dan periksa
        kesalinglepasannya.
  \item Empat pintu, satu mobil, tiga kambing. Pembawa acara yang tahu isi membuka dua
        pintu kambing di antara tiga pintu lain dan menawarkan perpindahan ke satu pintu
        tersisa. Hitung peluang menang bila berpindah, lalu periksa dengan simulasi
        $10^{4}$ permainan.
  \item Nilai klaim ini: ``Dua uji bebas bersensitivitas dan spesifisitas $0.99$
        memberi kepastian $0.9999$ bila keduanya positif, asalkan kelazimannya $0.1$
        persen.'' Hitung angka yang benar dan sebutkan asumsinya.
  \item Apakah independensi berpasangan untuk tiga kejadian cukup untuk menyimpulkan
        $P(A\cap B\cap C)=P(A)P(B)P(C)$? Bantah dengan contoh dari dua lemparan koin,
        dan hitung kedua nilainya.
\end{enumerate}

Petunjuk untuk tiga butir tersulit. Butir 2: pada ruang berhingga yang setiap hasilnya
berpeluang positif, $P(A\mid B)=0$ memaksa $A\cap B=\varnothing$, jadi pada lemparan dadu
klaim itu benar; contoh penyangkal harus datang dari ruang kontinu, misalnya
$A=\{1/2\}$ dan $B=(0,1)$ dengan peluang seragam, tempat $P(A)=0$ walaupun $A$ berisi
titik. Butir 7: mula-mula mobil di pintu pilihan berpeluang $1/4$; bila tidak, satu-satunya
pintu yang tidak dibuka pembawa acara pasti berisi mobil, jadi peluang berpindah $3/4$.
Butir 9: ambil $A$ koin pertama gambar, $B$ koin kedua gambar, dan $C$ kedua koin sama;
ketiga pasangan memberi $1/4$ sedangkan irisannya $1/4$, bukan $1/8$.
\end{latihan}

\begin{untukpeneliti}
Yang benar-benar terbukti di bab ini: ketiga aksioma beserta akibatnya; aturan Bayes
sebagai teorema yang diturunkan dari definisi peluang bersyarat; kemustahilan dua
kejadian saling lepas berpeluang positif untuk juga independen; dan ketergantungan
jawaban Monty Hall pada protokol, yang analisisnya dirumuskan tegas oleh Morgan,
Chaganty, Dahiya, dan Doviak, \emph{The American Statistician} 1991, pada artikel
``Let's Make a Deal: The Player's Dilemma''. Yang populer tetapi tidak terbukti di sini
adalah klaim bahwa format frekuensi alami memperbaiki penalaran kelaziman pada manusia;
itu klaim empiris tentang pembaca, bukan teorema, dan buktinya ada di literatur
psikologi pengambilan keputusan.
\end{untukpeneliti}

Sumber primer: Blitzstein dan Hwang, \emph{Introduction to Probability}, Bab 1 (Why
probability, Naive definition of probability, How to count, Story proofs, Non-naive
definition of probability) dan Bab 2 (The importance of thinking conditionally,
Definition and intuition, Bayes' rule and the law of total probability, Conditional
probabilities are probabilities, Independence of events, Coherency of Bayes' rule);
serta Wasserman, \emph{All of Statistics}, Bab 1 bagian Probability untuk perumusan
aksioma dan Bab 2 bagian Independent Events.

Batas keberlakuan: seluruh bab bekerja pada ruang berhingga yang setiap hasilnya
berpeluang positif, tempat $P(A)=0$ setara dengan $A=\varnothing$. Pada ruang tak
tercacah, pengkondisian pada kejadian berpeluang nol tidak didefinisikan oleh rumus
nisbah dan memerlukan distribusi bersyarat reguler; paradoks Borel menunjukkan jawabannya
dapat bergantung pada cara mengambil limit. Independensi juga tidak stabil: ia dapat
hilang atau muncul di bawah pengkondisian, dan pada model grafis muncul sebagai bias
seleksi.

Pertanyaan terbuka: seberapa jauh kelaziman harus diketahui sebelum aturan Bayes dipakai
secara jujur, dan apa yang dapat dikatakan bila kelaziman hanya diketahui dari sampel
kecil. Jawaban formalnya inferensi Bayesian dengan prior, yang membawa persoalan
pemilihan prior; jawaban frekuentisnya uji hipotesis di Bab 20, yang tidak memakai prior
tetapi menjawab pertanyaan berbeda. Bab ini tidak memilih salah satu.

\peta{Bab 17 pada Blitzstein dan Hwang, \emph{Introduction to Probability}: Bab 1
(Probability and counting, termasuk How to count dan contoh masalah ulang tahun) untuk
aksioma, pencacahan, dan peluang kombinatorial; Bab 2 (Conditional probability, Bayes'
rule and the law of total probability, Independence of events) untuk peluang bersyarat,
Bayes, dan independensi. Pada Wasserman, \emph{All of Statistics}: Bab 1 bagian
Probability dan Bab 2 bagian Independent Events. Himpunan dan komplemen dipakai dari
Bab 2 buku ini; deret yang mendasari aditif terbilang dari Bab 5; hukum bilangan besar
yang menopang simulasi dibuktikan di Bab 18. Variabel acak di Bab 17 memindahkan bobot
dari kejadian ke bilangan, sedangkan uji hipotesis di Bab 19 memakai aturan Bayes dan hukum
peluang total untuk pertanyaan berbeda.}